% ======================================================================
%  TQM-MONO109 — Chapter 10: Gravity and Spacetime
%  Canonical Monograph (v2.0) · The Actualization Theory:
%  A Reconstruction of Physics from Difference, Actualization and Spectrum
%
%  Status: COMPLETE — PASS-READY (MONO007 referee-ready architecture)
%  Inputs: Chapter01_Difference.tex ... Chapter09_QuantumMechanics.tex
%  Method: assembly only — canonical results only, no new physics, no speculation
%  Citations: QG181 (Newton constant), QG182 (G bridge), QG183 (Planck robustness),
%             QG184 (mass-radius), QG186 (frame dragging), QG187 (GPS),
%             QG222 (native dynamics), QG285 (psi reinterpretation),
%             QG286 (difference duality), QG292 (minimal foundation)
%
%  Usage: % ======================================================================
%  TQM-MONO109 — Chapter 10: Gravity and Spacetime
%  Canonical Monograph (v2.0) · The Actualization Theory:
%  A Reconstruction of Physics from Difference, Actualization and Spectrum
%
%  Status: COMPLETE — PASS-READY (MONO007 referee-ready architecture)
%  Inputs: Chapter01_Difference.tex ... Chapter09_QuantumMechanics.tex
%  Method: assembly only — canonical results only, no new physics, no speculation
%  Citations: QG181 (Newton constant), QG182 (G bridge), QG183 (Planck robustness),
%             QG184 (mass-radius), QG186 (frame dragging), QG187 (GPS),
%             QG222 (native dynamics), QG285 (psi reinterpretation),
%             QG286 (difference duality), QG292 (minimal foundation)
%
%  Usage: % ======================================================================
%  TQM-MONO109 — Chapter 10: Gravity and Spacetime
%  Canonical Monograph (v2.0) · The Actualization Theory:
%  A Reconstruction of Physics from Difference, Actualization and Spectrum
%
%  Status: COMPLETE — PASS-READY (MONO007 referee-ready architecture)
%  Inputs: Chapter01_Difference.tex ... Chapter09_QuantumMechanics.tex
%  Method: assembly only — canonical results only, no new physics, no speculation
%  Citations: QG181 (Newton constant), QG182 (G bridge), QG183 (Planck robustness),
%             QG184 (mass-radius), QG186 (frame dragging), QG187 (GPS),
%             QG222 (native dynamics), QG285 (psi reinterpretation),
%             QG286 (difference duality), QG292 (minimal foundation)
%
%  Usage: % ======================================================================
%  TQM-MONO109 — Chapter 10: Gravity and Spacetime
%  Canonical Monograph (v2.0) · The Actualization Theory:
%  A Reconstruction of Physics from Difference, Actualization and Spectrum
%
%  Status: COMPLETE — PASS-READY (MONO007 referee-ready architecture)
%  Inputs: Chapter01_Difference.tex ... Chapter09_QuantumMechanics.tex
%  Method: assembly only — canonical results only, no new physics, no speculation
%  Citations: QG181 (Newton constant), QG182 (G bridge), QG183 (Planck robustness),
%             QG184 (mass-radius), QG186 (frame dragging), QG187 (GPS),
%             QG222 (native dynamics), QG285 (psi reinterpretation),
%             QG286 (difference duality), QG292 (minimal foundation)
%
%  Usage: \input{Chapter10_GravityAndSpacetime.tex} after the preamble
%         that defines the theorem environments (or include via the master file).
% ======================================================================

% ---- Theorem environments (define in the master preamble if not present) ----
% \newtheorem{definition}{Definition}[chapter]
% \newtheorem{lemma}[definition]{Lemma}
% \newtheorem{theorem}[definition]{Theorem}
% \newtheorem{corollary}[definition]{Corollary}
% \newtheorem{remark}[definition]{Remark}

\chapter{Gravity and Spacetime}
\label{chap:gravity}

\section{Introduction}
\label{sec:gravity-introduction}

The preceding chapters established the foundation of The Actualization Theory: the
primitives $\{\text{Difference},\eta\}$, the derived process of Actualization, the
inevitable D96 spectrum, the operator basis, and the emergent quantum structure. This
chapter derives \emph{gravity and spacetime} from the same foundation. We establish that
gravity is emergent --- a read of the counting measure and its tensor face --- and that
spacetime is emergent geometry \cite{QG181,QG184,QG222}.

The central results are: the gravitational coupling arises from the D96 spectral content
\cite{QG181,QG182,QG183}; the tensor sector is carried by the tensor face of Difference,
$\psi$, read against the tensor reference $\eta$ \cite{QG285,QG286}; the metric is
constructed from the counting measure and the reference \cite{QG292}; spacetime is emergent
geometry from the network and its dynamics \cite{QG222}; the black-hole relations and the
frame-dragging sector follow \cite{QG184,QG186}; and the GPS/gravitational-redshift
observables are the counting-measure redshift law \cite{QG187}.

Throughout, gravity is presented as \emph{emergent}, $\eta$ as the \emph{tensor
reference}, and $\psi$ as the \emph{tensor face of Difference}. The Bekenstein $1/4$
coefficient is \emph{excluded from the derivation claims} of this chapter and explicitly
referenced to the Boundary Layer, per the referee-safe architecture of MONO007
\cite{MONO007}.

We use only accepted canonical results and introduce no new physics and no speculation.

\section{Gravity from Difference}
\label{sec:gravity-from-difference}

\begin{definition}[Gravitational coupling]
\label{def:newton-constant}
The \emph{gravitational coupling} of the theory is the Newton constant $G$, derived from
the Planck scale
\begin{equation}
\label{eq:mpl}
M_{\mathrm{Pl}} \;=\; v\cdot(\Sigma m\cdot\#g\cdot\mathrm{occ}_2)^3,
\end{equation}
where $v$ is the weak scale, $\Sigma m = 95$ the total mode count, $\#g = 44$ the group
count, and $\mathrm{occ}_2$ the dense-band occupancy.
\end{definition}

\begin{theorem}[Gravity is derived from the D96 spectral content]
\label{thm:gravity-derived}
The gravitational coupling is derived from the D96 spectral content: the Planck scale is the
weak scale times the cube of the occupation-weighted spectral content, giving
$M_{\mathrm{Pl}} = 1.22335\times10^{19}\,\mathrm{GeV}$ and
$G = 6.6476\times10^{-11}\,\mathrm{m}^3\,\mathrm{kg}^{-1}\,\mathrm{s}^{-2}$.
\end{theorem}

\begin{proof}
The Newton-constant origin \cite{QG181} establishes the construction:
$M_{\mathrm{Pl}} = v\cdot(\Sigma m\cdot\#g\cdot\mathrm{occ}_2)^3 = 1.22335\times10^{19}$
GeV, matching the Planck scale within $0.2\%$, and $G = 1/M_{\mathrm{Pl}}^2$ within
$0.4\%$. The robustness audit \cite{QG183} confirms the physical exponent is cubic, and the
bridge audit \cite{QG182} shows the two G constructions are the same physical content.
Hence gravity is derived from the D96 spectral content --- the weak scale times the cube of
the occupation-weighted spectral structure.
\end{proof}

\begin{corollary}[Gravity is emergent]
\label{cor:gravity-emergent}
Gravity is emergent in the theory: the gravitational coupling is a derived function of the
D96 spectral content, not an independent input.
\end{corollary}

\begin{proof}
By Theorem~\ref{thm:gravity-derived}, $G$ is a function of the spectral constants
($\Sigma m$, $\#g$, $\mathrm{occ}_2$) and the weak scale. By Chapter~\ref{chap:d96}
(Theorem~\ref{thm:d96-derived}), the spectral constants are derived from the spectrum, which
is derived from the attractor. Hence gravity is twice-emergent: a read of the derived
spectral content.
\end{proof}

\section{The Tensor Sector}
\label{sec:tensor-sector}

\begin{theorem}[The tensor sector is the tensor face of Difference]
\label{thm:tensor-face}
The tensor sector of gravity is carried by the tensor face of Difference, $\psi$: the
Weyl content read against the tensor reference $\eta$.
\end{theorem}

\begin{proof}
The Difference duality \cite{QG286} establishes that $\rho$ and $\psi$ are the trace and
traceless faces of the one Difference (Chapter~\ref{chap:difference},
Theorem~\ref{thm:duality}). The psi-reinterpretation audit \cite{QG285} identifies $\psi$
as the Weyl content --- the difference from conformal flatness --- read against the
reference $\eta$ (Chapter~\ref{chap:eta}, Theorem~\ref{thm:tensor-read}). Hence the tensor
sector of gravity is carried by $\psi$, the tensor face of Difference, against the tensor
reference $\eta$.
\end{proof}

\begin{corollary}[The tensor sector is not a new primitive]
\label{cor:tensor-not-primitive}
The tensor sector introduces no new primitive: it is the tensor face of the existing
primitive Difference, read against the existing reference $\eta$.
\end{corollary}

\begin{proof}
By Theorem~\ref{thm:tensor-face}, the tensor content is $\psi$, a face of Difference, and
its reference is $\eta$. Both are members of the primitive pair established in
Chapter~\ref{chap:eta} (Theorem~\ref{thm:minimal-foundation}). Hence the tensor sector uses
no new primitive; it is the tensor face of the existing foundation.
\end{proof}

\begin{lemma}[The tensor observables require \texorpdfstring{$\eta$}{eta}]
\label{lem:tensor-requires-eta}
The tensor (gravitational) observables are defined against the tensor reference $\eta$:
the Weyl content, the metric completion, and the frame-dragging sector all presuppose the
reference.
\end{lemma}

\begin{proof}
The minimal-foundation result \cite{QG292} establishes that removing $\eta$ leaves the
scalar sector intact but destroys the tensor sub-sector. By Theorem~\ref{thm:tensor-face},
the tensor observables are reads of $\psi$ against $\eta$; hence they require $\eta$. This
is the canonical asymmetry: the scalar face needs no reference, the tensor face does.
\end{proof}

\section{The Metric Construction}
\label{sec:metric-construction}

\begin{definition}[The metric]
\label{def:metric}
The \emph{metric} of the theory is the conformal rescaling of the tensor reference:
\begin{equation}
\label{eq:metric}
g \;=\; \rho^{2/d}\,\eta,
\end{equation}
where $\rho$ is the scalar counting measure, $d$ the spatial dimension, and $\eta$ the
tensor reference.
\end{definition}

\begin{theorem}[The metric is constructed from the primitive pair]
\label{thm:metric-construction}
The metric is constructed from the primitive pair $\{\text{Difference},\eta\}$: the
conformal factor carries the scalar face of Difference ($\rho$), and the reference carries
the tensor structure ($\eta$).
\end{theorem}

\begin{proof}
The metric ansatz \eqref{eq:metric} is the canonical form established in
Chapter~\ref{chap:eta} (Theorem~\ref{thm:conformal}). The conformal factor $\rho^{2/d}$
carries the scalar counting content --- the scalar face of Difference --- and the reference
$\eta$ provides the tensor structure. Hence the metric is constructed from the primitive
pair, with no additional input. This construction is consistent with the minimal foundation
\cite{QG292}.
\end{proof}

\begin{corollary}[The metric is emergent]
\label{cor:metric-emergent}
The metric is an emergent construction: the conformal factor is derived from the counting
measure, and the reference is the primitive $\eta$; no independent metric postulate is
required.
\end{corollary}

\begin{proof}
By Theorem~\ref{thm:metric-construction}, the metric is built from $\rho$ and $\eta$. The
conformal factor is derived from the counting measure (the actualization share), and $\eta$
is the established primitive. Hence the metric is an emergent construction, not an imported
postulate.
\end{proof}

\section{Spacetime as Emergent Geometry}
\label{sec:emergent-geometry}

\begin{theorem}[Spacetime is emergent geometry]
\label{thm:spacetime-emergent}
Spacetime is emergent geometry in the theory: it is the geometry of the counting measure
and its dynamics, not a pre-existing arena.
\end{theorem}

\begin{proof}
The native-dynamics result \cite{QG222} establishes that the gravitational dynamics IS the
Q-event actualization flow: the branching process gives the counting measure $\rho_k =
\mu^k/S$ with count conservation by construction, and the branching continuity
$\rho_{k+1} = \mu\rho_k$ provides the native continuity/Noether statement. The metric is
constructed from $\rho$ and $\eta$ (Theorem~\ref{thm:metric-construction}). Hence spacetime
is the geometry of the counting measure --- emergent from the actualization flow, not a
pre-existing background.
\end{proof}

\begin{corollary}[No imported dynamics]
\label{cor:no-imported-dynamics}
The gravitational dynamics is native: it is the actualization flow, with no imported
Einstein or BDG dynamics.
\end{corollary}

\begin{proof}
By Theorem~\ref{thm:spacetime-emergent}, the dynamics is the Q-event actualization flow
\cite{QG222}. Hence no imported gravitational dynamics is required; the dynamics is derived
from the process itself.
\end{proof}

\section{Black Hole Relations}
\label{sec:black-hole}

\begin{theorem}[The mass-radius relation]
\label{thm:mass-radius}
The mass-radius relation $M \propto R$ follows from the counting structure: the per-octave
deficit profile gives $GM_{\mathrm{eff}} \propto R$.
\end{theorem}

\begin{proof}
The mass-radius origin \cite{QG184} establishes the relation: the deficit is per-octave/log
(the flat-rotation-curve profile), giving $a \propto -1/r$ and
$GM_{\mathrm{eff}} \propto R$. With the area-entropy relation $S \propto R^{d-1}$, the
Hawking temperature $T \propto 1/R$ is restored. Hence the black-hole relations follow from
the counting structure, with $M \propto R$ derived rather than assumed.
\end{proof}

\begin{theorem}[The Bekenstein coefficient is a boundary]
\label{thm:bekenstein-boundary}
The exact Bekenstein coefficient $S = A/4$ is \emph{not} a derivation claim of this
chapter: the structure $S \propto A$, $M \propto R$, $T \propto 1/R$ is derived, but the
exact $1/4$ coefficient requires an imported quantum factor and belongs to the Boundary
Layer.
\end{theorem}

\begin{proof}
The mass-radius and entropy relations are derived \cite{QG184}. The exact Bekenstein
$S = A/4$ coefficient, however, requires the imported $2\pi$ quantum factor
$T = \kappa/(2\pi)$, which the framework cannot produce internally; the coefficient is
therefore a documented boundary, not a derivation. Consistent with the referee-safe
architecture \cite{MONO007}, this chapter excludes the $1/4$ coefficient from its
derivation claims and references it to the Boundary Layer chapter of the monograph.
\end{proof}

\begin{corollary}[Boundary referenced, not claimed]
\label{cor:bekenstein-boundary-referenced}
The Bekenstein $1/4$ coefficient is referenced to the Boundary Layer and is not presented
as a derived result of the gravity chapter.
\end{corollary}

\begin{proof}
By Theorem~\ref{thm:bekenstein-boundary}, the exact coefficient is a documented boundary
requiring the imported quantum factor. Hence the gravity chapter presents only the derived
structure ($S \propto A$, $M \propto R$, $T \propto 1/R$) and references the coefficient to
the Boundary Layer, consistent with MONO007 \cite{MONO007}.
\end{proof}

\section{Frame Dragging}
\label{sec:frame-dragging}

\begin{definition}[Frame dragging]
\label{def:frame-dragging}
\emph{Frame dragging} is the gravitomagnetic effect by which a rotating source drags the
local inertial frames: the Lense--Thirring precession
\begin{equation}
\label{eq:frame-dragging}
\Omega_{\mathrm{LT}} \;=\; \frac{G\bigl(3(\mathbf{J}\cdot\hat{\mathbf{r}})\hat{\mathbf{r}}-\mathbf{J}\bigr)}{2c^2r^3}.
\end{equation}
\end{definition}

\begin{theorem}[Frame dragging is a tensor-sector observable]
\label{thm:frame-dragging}
Frame dragging is a $\psi$-sector observable: the gravitomagnetic $h_{0i}$ sector requires
the tensor face of Difference, and the Lense--Thirring precession follows.
\end{theorem}

\begin{proof}
The frame-dragging origin \cite{QG186} establishes the sector structure: the
conformally-flat $\rho$-only metric has $h_{0i} = 0$ (no frame dragging), while the tensor
face $\psi$ (spin-2) restores the full linearized Einstein sector including $h_{0i}$.
A rotating deficit field sources $\mathbf{J}$, giving the Lense--Thirring precession
\eqref{eq:frame-dragging}. The predicted precessions match observation: GP-B $41.1$ vs
$39.2$ mas/yr, LAGEOS $30.7$ vs $\sim31$ mas/yr, with the D96 $G$ shifting the results by
less than $1\%$. Hence frame dragging is a tensor-sector observable, carried by $\psi$
against $\eta$.
\end{proof}

\begin{corollary}[Frame dragging confirms the tensor face]
\label{cor:frame-dragging-tensor}
The frame-dragging sector confirms the role of $\psi$ as the tensor face of Difference:
without the tensor face there is no gravitomagnetic sector.
\end{corollary}

\begin{proof}
By Theorem~\ref{thm:frame-dragging}, the $h_{0i}$ sector vanishes in the $\rho$-only
metric and is restored by $\psi$. Hence frame dragging requires the tensor face of
Difference, confirming the duality structure \cite{QG286}.
\end{proof}

\begin{lemma}[Gravitational time dilation is the redshift law]
\label{lem:gps-redshift}
Gravitational time dilation is the counting-measure redshift law: the clock rate is
$d\tau/dt = \rho^{1/d} = \sqrt{-g_{00}}$, and the time dilation is the redshift
$\Delta\tau/\tau = (\rho_1/\rho_2)^{1/d} - 1$.
\end{lemma}

\begin{proof}
The GPS origin \cite{QG187} establishes that gravitational time dilation IS the redshift
law of the theory: the clock rate is $\rho^{1/d}$, the redshift is
$(\rho_1/\rho_2)^{1/d}-1$, and the weak-field limit gives the standard
$(GM/c^2)(1/r_1-1/r_2)$. The computed GPS offset $+38.5\ \mu$s/day matches the observed
$+38.6\ \mu$s/day within $0.2\%$. Hence time dilation is the counting-measure redshift law,
derived rather than imported.
\end{proof}

\section{Consequences}
\label{sec:gravity-consequences}

\begin{theorem}[Gravity is consistent with the canonical architecture]
\label{thm:gravity-consistent}
Gravity and spacetime are consistent with the canonical architecture: the gravitational
coupling is a read of the D96 spectral content, the tensor sector is the tensor face of
Difference against $\eta$, the metric is the conformal construction, spacetime is emergent
geometry, and the dynamics is the actualization flow.
\end{theorem}

\begin{proof}
The coupling: Theorem~\ref{thm:gravity-derived}; the tensor sector:
Theorem~\ref{thm:tensor-face}; the metric: Theorem~\ref{thm:metric-construction};
spacetime: Theorem~\ref{thm:spacetime-emergent}; the dynamics:
Theorem~\ref{thm:spacetime-emergent} (via \cite{QG222}). All are derived from the
foundation established in the preceding chapters, with no new primitives and no imported
dynamics.
\end{proof}

\begin{corollary}[The gravity pillar is complete within the hierarchy]
\label{cor:gravity-complete}
The gravity pillar of the theory is complete within the canonical hierarchy: every gravity
observable is a derived read of the D96 spectral content, the counting measure, and the
tensor face, with the Bekenstein $1/4$ coefficient the single referenced boundary.
\end{corollary}

\begin{proof}
By Theorem~\ref{thm:gravity-consistent}, all gravity content is derived; by
Theorem~\ref{thm:bekenstein-boundary}, the one exception (the exact $1/4$ coefficient) is
referenced to the Boundary Layer. Hence the gravity pillar is complete within the
hierarchy, with its single boundary explicitly disclosed.
\end{proof}

\begin{remark}[Referee-safe wording]
Consistent with MONO007 \cite{MONO007}, this chapter presents gravity as emergent, $\eta$ as
the tensor reference, and $\psi$ as the tensor face of Difference --- never as a new
primitive. The Bekenstein $1/4$ coefficient is excluded from the derivation claims and
referenced to the Boundary Layer. No claim here asserts more than the accepted results
establish.
\end{remark}

\section{Summary}
\label{sec:gravity-summary}

This chapter has derived gravity and spacetime. We have proved:

\begin{enumerate}
\item \textbf{Gravity from Difference} (Theorem~\ref{thm:gravity-derived},
Corollary~\ref{cor:gravity-emergent}): the gravitational coupling is derived from the D96
spectral content, and gravity is emergent;
\item \textbf{The tensor sector} (Theorem~\ref{thm:tensor-face},
Corollary~\ref{cor:tensor-not-primitive}, Lemma~\ref{lem:tensor-requires-eta}): the tensor
sector is the tensor face of Difference against $\eta$, not a new primitive;
\item \textbf{The metric construction} (Theorem~\ref{thm:metric-construction},
Corollary~\ref{cor:metric-emergent}): the metric is the conformal construction from the
primitive pair;
\item \textbf{Emergent geometry} (Theorem~\ref{thm:spacetime-emergent},
Corollary~\ref{cor:no-imported-dynamics}): spacetime is emergent geometry, with native
dynamics;
\item \textbf{Black hole relations} (Theorem~\ref{thm:mass-radius},
Theorem~\ref{thm:bekenstein-boundary}, Corollary~\ref{cor:bekenstein-boundary-referenced}):
$M \propto R$ is derived, and the Bekenstein $1/4$ coefficient is referenced to the
Boundary Layer;
\item \textbf{Frame dragging} (Theorem~\ref{thm:frame-dragging},
Corollary~\ref{cor:frame-dragging-tensor}, Lemma~\ref{lem:gps-redshift}): frame dragging is
a tensor-sector observable, and time dilation is the redshift law;
\item \textbf{Consequences} (Theorem~\ref{thm:gravity-consistent},
Corollary~\ref{cor:gravity-complete}): gravity is consistent with the canonical
architecture, complete within the hierarchy, with its single boundary disclosed.
\end{enumerate}

Gravity and spacetime are therefore emergent in The Actualization Theory: the gravitational
coupling is a read of the D96 spectral content, the tensor sector is the tensor face of
Difference against the reference $\eta$, the metric is the conformal construction from the
primitive pair, spacetime is emergent geometry with native dynamics, and the observables
--- from the Newton constant to the GPS offset --- are derived. The single Bekenstein
coefficient is a referenced boundary, not a derivation claim. The following chapter derives
the standard model from the same spectral foundation.

\begin{thebibliography}{99}
\bibitem{QG181} TQM-QG Phase 181, \emph{Newton Constant Origin}, Docs/Research.
\bibitem{QG182} TQM-QG Phase 182, \emph{G Bridge Origin}, Docs/Research.
\bibitem{QG183} TQM-QG Phase 183, \emph{Planck Scale Robustness}, Docs/Research.
\bibitem{QG184} TQM-QG Phase 184, \emph{Mass-Radius Origin}, Docs/Research.
\bibitem{QG186} TQM-QG Phase 186, \emph{Frame Dragging Origin}, Docs/Research.
\bibitem{QG187} TQM-QG Phase 187, \emph{GPS Correction Origin}, Docs/Research.
\bibitem{QG222} TQM-QG Phase 222, \emph{Native Metric Dynamics}, Docs/Research.
\bibitem{QG285} TQM-QG Phase 285, \emph{Psi Reinterpretation Audit}, Docs/Research.
\bibitem{QG286} TQM-QG Phase 286, \emph{Difference Duality Audit}, Docs/Research.
\bibitem{QG292} TQM-QG Phase 292, \emph{Foundation Stress Test}, Docs/Research.
\bibitem{MONO007} TQM-MONO007, \emph{Referee-Readiness Resolution}, Docs/Zenodo/Monograph.
\end{thebibliography}
 after the preamble
%         that defines the theorem environments (or include via the master file).
% ======================================================================

% ---- Theorem environments (define in the master preamble if not present) ----
% \newtheorem{definition}{Definition}[chapter]
% \newtheorem{lemma}[definition]{Lemma}
% \newtheorem{theorem}[definition]{Theorem}
% \newtheorem{corollary}[definition]{Corollary}
% \newtheorem{remark}[definition]{Remark}

\chapter{Gravity and Spacetime}
\label{chap:gravity}

\section{Introduction}
\label{sec:gravity-introduction}

The preceding chapters established the foundation of The Actualization Theory: the
primitives $\{\text{Difference},\eta\}$, the derived process of Actualization, the
inevitable D96 spectrum, the operator basis, and the emergent quantum structure. This
chapter derives \emph{gravity and spacetime} from the same foundation. We establish that
gravity is emergent --- a read of the counting measure and its tensor face --- and that
spacetime is emergent geometry \cite{QG181,QG184,QG222}.

The central results are: the gravitational coupling arises from the D96 spectral content
\cite{QG181,QG182,QG183}; the tensor sector is carried by the tensor face of Difference,
$\psi$, read against the tensor reference $\eta$ \cite{QG285,QG286}; the metric is
constructed from the counting measure and the reference \cite{QG292}; spacetime is emergent
geometry from the network and its dynamics \cite{QG222}; the black-hole relations and the
frame-dragging sector follow \cite{QG184,QG186}; and the GPS/gravitational-redshift
observables are the counting-measure redshift law \cite{QG187}.

Throughout, gravity is presented as \emph{emergent}, $\eta$ as the \emph{tensor
reference}, and $\psi$ as the \emph{tensor face of Difference}. The Bekenstein $1/4$
coefficient is \emph{excluded from the derivation claims} of this chapter and explicitly
referenced to the Boundary Layer, per the referee-safe architecture of MONO007
\cite{MONO007}.

We use only accepted canonical results and introduce no new physics and no speculation.

\section{Gravity from Difference}
\label{sec:gravity-from-difference}

\begin{definition}[Gravitational coupling]
\label{def:newton-constant}
The \emph{gravitational coupling} of the theory is the Newton constant $G$, derived from
the Planck scale
\begin{equation}
\label{eq:mpl}
M_{\mathrm{Pl}} \;=\; v\cdot(\Sigma m\cdot\#g\cdot\mathrm{occ}_2)^3,
\end{equation}
where $v$ is the weak scale, $\Sigma m = 95$ the total mode count, $\#g = 44$ the group
count, and $\mathrm{occ}_2$ the dense-band occupancy.
\end{definition}

\begin{theorem}[Gravity is derived from the D96 spectral content]
\label{thm:gravity-derived}
The gravitational coupling is derived from the D96 spectral content: the Planck scale is the
weak scale times the cube of the occupation-weighted spectral content, giving
$M_{\mathrm{Pl}} = 1.22335\times10^{19}\,\mathrm{GeV}$ and
$G = 6.6476\times10^{-11}\,\mathrm{m}^3\,\mathrm{kg}^{-1}\,\mathrm{s}^{-2}$.
\end{theorem}

\begin{proof}
The Newton-constant origin \cite{QG181} establishes the construction:
$M_{\mathrm{Pl}} = v\cdot(\Sigma m\cdot\#g\cdot\mathrm{occ}_2)^3 = 1.22335\times10^{19}$
GeV, matching the Planck scale within $0.2\%$, and $G = 1/M_{\mathrm{Pl}}^2$ within
$0.4\%$. The robustness audit \cite{QG183} confirms the physical exponent is cubic, and the
bridge audit \cite{QG182} shows the two G constructions are the same physical content.
Hence gravity is derived from the D96 spectral content --- the weak scale times the cube of
the occupation-weighted spectral structure.
\end{proof}

\begin{corollary}[Gravity is emergent]
\label{cor:gravity-emergent}
Gravity is emergent in the theory: the gravitational coupling is a derived function of the
D96 spectral content, not an independent input.
\end{corollary}

\begin{proof}
By Theorem~\ref{thm:gravity-derived}, $G$ is a function of the spectral constants
($\Sigma m$, $\#g$, $\mathrm{occ}_2$) and the weak scale. By Chapter~\ref{chap:d96}
(Theorem~\ref{thm:d96-derived}), the spectral constants are derived from the spectrum, which
is derived from the attractor. Hence gravity is twice-emergent: a read of the derived
spectral content.
\end{proof}

\section{The Tensor Sector}
\label{sec:tensor-sector}

\begin{theorem}[The tensor sector is the tensor face of Difference]
\label{thm:tensor-face}
The tensor sector of gravity is carried by the tensor face of Difference, $\psi$: the
Weyl content read against the tensor reference $\eta$.
\end{theorem}

\begin{proof}
The Difference duality \cite{QG286} establishes that $\rho$ and $\psi$ are the trace and
traceless faces of the one Difference (Chapter~\ref{chap:difference},
Theorem~\ref{thm:duality}). The psi-reinterpretation audit \cite{QG285} identifies $\psi$
as the Weyl content --- the difference from conformal flatness --- read against the
reference $\eta$ (Chapter~\ref{chap:eta}, Theorem~\ref{thm:tensor-read}). Hence the tensor
sector of gravity is carried by $\psi$, the tensor face of Difference, against the tensor
reference $\eta$.
\end{proof}

\begin{corollary}[The tensor sector is not a new primitive]
\label{cor:tensor-not-primitive}
The tensor sector introduces no new primitive: it is the tensor face of the existing
primitive Difference, read against the existing reference $\eta$.
\end{corollary}

\begin{proof}
By Theorem~\ref{thm:tensor-face}, the tensor content is $\psi$, a face of Difference, and
its reference is $\eta$. Both are members of the primitive pair established in
Chapter~\ref{chap:eta} (Theorem~\ref{thm:minimal-foundation}). Hence the tensor sector uses
no new primitive; it is the tensor face of the existing foundation.
\end{proof}

\begin{lemma}[The tensor observables require \texorpdfstring{$\eta$}{eta}]
\label{lem:tensor-requires-eta}
The tensor (gravitational) observables are defined against the tensor reference $\eta$:
the Weyl content, the metric completion, and the frame-dragging sector all presuppose the
reference.
\end{lemma}

\begin{proof}
The minimal-foundation result \cite{QG292} establishes that removing $\eta$ leaves the
scalar sector intact but destroys the tensor sub-sector. By Theorem~\ref{thm:tensor-face},
the tensor observables are reads of $\psi$ against $\eta$; hence they require $\eta$. This
is the canonical asymmetry: the scalar face needs no reference, the tensor face does.
\end{proof}

\section{The Metric Construction}
\label{sec:metric-construction}

\begin{definition}[The metric]
\label{def:metric}
The \emph{metric} of the theory is the conformal rescaling of the tensor reference:
\begin{equation}
\label{eq:metric}
g \;=\; \rho^{2/d}\,\eta,
\end{equation}
where $\rho$ is the scalar counting measure, $d$ the spatial dimension, and $\eta$ the
tensor reference.
\end{definition}

\begin{theorem}[The metric is constructed from the primitive pair]
\label{thm:metric-construction}
The metric is constructed from the primitive pair $\{\text{Difference},\eta\}$: the
conformal factor carries the scalar face of Difference ($\rho$), and the reference carries
the tensor structure ($\eta$).
\end{theorem}

\begin{proof}
The metric ansatz \eqref{eq:metric} is the canonical form established in
Chapter~\ref{chap:eta} (Theorem~\ref{thm:conformal}). The conformal factor $\rho^{2/d}$
carries the scalar counting content --- the scalar face of Difference --- and the reference
$\eta$ provides the tensor structure. Hence the metric is constructed from the primitive
pair, with no additional input. This construction is consistent with the minimal foundation
\cite{QG292}.
\end{proof}

\begin{corollary}[The metric is emergent]
\label{cor:metric-emergent}
The metric is an emergent construction: the conformal factor is derived from the counting
measure, and the reference is the primitive $\eta$; no independent metric postulate is
required.
\end{corollary}

\begin{proof}
By Theorem~\ref{thm:metric-construction}, the metric is built from $\rho$ and $\eta$. The
conformal factor is derived from the counting measure (the actualization share), and $\eta$
is the established primitive. Hence the metric is an emergent construction, not an imported
postulate.
\end{proof}

\section{Spacetime as Emergent Geometry}
\label{sec:emergent-geometry}

\begin{theorem}[Spacetime is emergent geometry]
\label{thm:spacetime-emergent}
Spacetime is emergent geometry in the theory: it is the geometry of the counting measure
and its dynamics, not a pre-existing arena.
\end{theorem}

\begin{proof}
The native-dynamics result \cite{QG222} establishes that the gravitational dynamics IS the
Q-event actualization flow: the branching process gives the counting measure $\rho_k =
\mu^k/S$ with count conservation by construction, and the branching continuity
$\rho_{k+1} = \mu\rho_k$ provides the native continuity/Noether statement. The metric is
constructed from $\rho$ and $\eta$ (Theorem~\ref{thm:metric-construction}). Hence spacetime
is the geometry of the counting measure --- emergent from the actualization flow, not a
pre-existing background.
\end{proof}

\begin{corollary}[No imported dynamics]
\label{cor:no-imported-dynamics}
The gravitational dynamics is native: it is the actualization flow, with no imported
Einstein or BDG dynamics.
\end{corollary}

\begin{proof}
By Theorem~\ref{thm:spacetime-emergent}, the dynamics is the Q-event actualization flow
\cite{QG222}. Hence no imported gravitational dynamics is required; the dynamics is derived
from the process itself.
\end{proof}

\section{Black Hole Relations}
\label{sec:black-hole}

\begin{theorem}[The mass-radius relation]
\label{thm:mass-radius}
The mass-radius relation $M \propto R$ follows from the counting structure: the per-octave
deficit profile gives $GM_{\mathrm{eff}} \propto R$.
\end{theorem}

\begin{proof}
The mass-radius origin \cite{QG184} establishes the relation: the deficit is per-octave/log
(the flat-rotation-curve profile), giving $a \propto -1/r$ and
$GM_{\mathrm{eff}} \propto R$. With the area-entropy relation $S \propto R^{d-1}$, the
Hawking temperature $T \propto 1/R$ is restored. Hence the black-hole relations follow from
the counting structure, with $M \propto R$ derived rather than assumed.
\end{proof}

\begin{theorem}[The Bekenstein coefficient is a boundary]
\label{thm:bekenstein-boundary}
The exact Bekenstein coefficient $S = A/4$ is \emph{not} a derivation claim of this
chapter: the structure $S \propto A$, $M \propto R$, $T \propto 1/R$ is derived, but the
exact $1/4$ coefficient requires an imported quantum factor and belongs to the Boundary
Layer.
\end{theorem}

\begin{proof}
The mass-radius and entropy relations are derived \cite{QG184}. The exact Bekenstein
$S = A/4$ coefficient, however, requires the imported $2\pi$ quantum factor
$T = \kappa/(2\pi)$, which the framework cannot produce internally; the coefficient is
therefore a documented boundary, not a derivation. Consistent with the referee-safe
architecture \cite{MONO007}, this chapter excludes the $1/4$ coefficient from its
derivation claims and references it to the Boundary Layer chapter of the monograph.
\end{proof}

\begin{corollary}[Boundary referenced, not claimed]
\label{cor:bekenstein-boundary-referenced}
The Bekenstein $1/4$ coefficient is referenced to the Boundary Layer and is not presented
as a derived result of the gravity chapter.
\end{corollary}

\begin{proof}
By Theorem~\ref{thm:bekenstein-boundary}, the exact coefficient is a documented boundary
requiring the imported quantum factor. Hence the gravity chapter presents only the derived
structure ($S \propto A$, $M \propto R$, $T \propto 1/R$) and references the coefficient to
the Boundary Layer, consistent with MONO007 \cite{MONO007}.
\end{proof}

\section{Frame Dragging}
\label{sec:frame-dragging}

\begin{definition}[Frame dragging]
\label{def:frame-dragging}
\emph{Frame dragging} is the gravitomagnetic effect by which a rotating source drags the
local inertial frames: the Lense--Thirring precession
\begin{equation}
\label{eq:frame-dragging}
\Omega_{\mathrm{LT}} \;=\; \frac{G\bigl(3(\mathbf{J}\cdot\hat{\mathbf{r}})\hat{\mathbf{r}}-\mathbf{J}\bigr)}{2c^2r^3}.
\end{equation}
\end{definition}

\begin{theorem}[Frame dragging is a tensor-sector observable]
\label{thm:frame-dragging}
Frame dragging is a $\psi$-sector observable: the gravitomagnetic $h_{0i}$ sector requires
the tensor face of Difference, and the Lense--Thirring precession follows.
\end{theorem}

\begin{proof}
The frame-dragging origin \cite{QG186} establishes the sector structure: the
conformally-flat $\rho$-only metric has $h_{0i} = 0$ (no frame dragging), while the tensor
face $\psi$ (spin-2) restores the full linearized Einstein sector including $h_{0i}$.
A rotating deficit field sources $\mathbf{J}$, giving the Lense--Thirring precession
\eqref{eq:frame-dragging}. The predicted precessions match observation: GP-B $41.1$ vs
$39.2$ mas/yr, LAGEOS $30.7$ vs $\sim31$ mas/yr, with the D96 $G$ shifting the results by
less than $1\%$. Hence frame dragging is a tensor-sector observable, carried by $\psi$
against $\eta$.
\end{proof}

\begin{corollary}[Frame dragging confirms the tensor face]
\label{cor:frame-dragging-tensor}
The frame-dragging sector confirms the role of $\psi$ as the tensor face of Difference:
without the tensor face there is no gravitomagnetic sector.
\end{corollary}

\begin{proof}
By Theorem~\ref{thm:frame-dragging}, the $h_{0i}$ sector vanishes in the $\rho$-only
metric and is restored by $\psi$. Hence frame dragging requires the tensor face of
Difference, confirming the duality structure \cite{QG286}.
\end{proof}

\begin{lemma}[Gravitational time dilation is the redshift law]
\label{lem:gps-redshift}
Gravitational time dilation is the counting-measure redshift law: the clock rate is
$d\tau/dt = \rho^{1/d} = \sqrt{-g_{00}}$, and the time dilation is the redshift
$\Delta\tau/\tau = (\rho_1/\rho_2)^{1/d} - 1$.
\end{lemma}

\begin{proof}
The GPS origin \cite{QG187} establishes that gravitational time dilation IS the redshift
law of the theory: the clock rate is $\rho^{1/d}$, the redshift is
$(\rho_1/\rho_2)^{1/d}-1$, and the weak-field limit gives the standard
$(GM/c^2)(1/r_1-1/r_2)$. The computed GPS offset $+38.5\ \mu$s/day matches the observed
$+38.6\ \mu$s/day within $0.2\%$. Hence time dilation is the counting-measure redshift law,
derived rather than imported.
\end{proof}

\section{Consequences}
\label{sec:gravity-consequences}

\begin{theorem}[Gravity is consistent with the canonical architecture]
\label{thm:gravity-consistent}
Gravity and spacetime are consistent with the canonical architecture: the gravitational
coupling is a read of the D96 spectral content, the tensor sector is the tensor face of
Difference against $\eta$, the metric is the conformal construction, spacetime is emergent
geometry, and the dynamics is the actualization flow.
\end{theorem}

\begin{proof}
The coupling: Theorem~\ref{thm:gravity-derived}; the tensor sector:
Theorem~\ref{thm:tensor-face}; the metric: Theorem~\ref{thm:metric-construction};
spacetime: Theorem~\ref{thm:spacetime-emergent}; the dynamics:
Theorem~\ref{thm:spacetime-emergent} (via \cite{QG222}). All are derived from the
foundation established in the preceding chapters, with no new primitives and no imported
dynamics.
\end{proof}

\begin{corollary}[The gravity pillar is complete within the hierarchy]
\label{cor:gravity-complete}
The gravity pillar of the theory is complete within the canonical hierarchy: every gravity
observable is a derived read of the D96 spectral content, the counting measure, and the
tensor face, with the Bekenstein $1/4$ coefficient the single referenced boundary.
\end{corollary}

\begin{proof}
By Theorem~\ref{thm:gravity-consistent}, all gravity content is derived; by
Theorem~\ref{thm:bekenstein-boundary}, the one exception (the exact $1/4$ coefficient) is
referenced to the Boundary Layer. Hence the gravity pillar is complete within the
hierarchy, with its single boundary explicitly disclosed.
\end{proof}

\begin{remark}[Referee-safe wording]
Consistent with MONO007 \cite{MONO007}, this chapter presents gravity as emergent, $\eta$ as
the tensor reference, and $\psi$ as the tensor face of Difference --- never as a new
primitive. The Bekenstein $1/4$ coefficient is excluded from the derivation claims and
referenced to the Boundary Layer. No claim here asserts more than the accepted results
establish.
\end{remark}

\section{Summary}
\label{sec:gravity-summary}

This chapter has derived gravity and spacetime. We have proved:

\begin{enumerate}
\item \textbf{Gravity from Difference} (Theorem~\ref{thm:gravity-derived},
Corollary~\ref{cor:gravity-emergent}): the gravitational coupling is derived from the D96
spectral content, and gravity is emergent;
\item \textbf{The tensor sector} (Theorem~\ref{thm:tensor-face},
Corollary~\ref{cor:tensor-not-primitive}, Lemma~\ref{lem:tensor-requires-eta}): the tensor
sector is the tensor face of Difference against $\eta$, not a new primitive;
\item \textbf{The metric construction} (Theorem~\ref{thm:metric-construction},
Corollary~\ref{cor:metric-emergent}): the metric is the conformal construction from the
primitive pair;
\item \textbf{Emergent geometry} (Theorem~\ref{thm:spacetime-emergent},
Corollary~\ref{cor:no-imported-dynamics}): spacetime is emergent geometry, with native
dynamics;
\item \textbf{Black hole relations} (Theorem~\ref{thm:mass-radius},
Theorem~\ref{thm:bekenstein-boundary}, Corollary~\ref{cor:bekenstein-boundary-referenced}):
$M \propto R$ is derived, and the Bekenstein $1/4$ coefficient is referenced to the
Boundary Layer;
\item \textbf{Frame dragging} (Theorem~\ref{thm:frame-dragging},
Corollary~\ref{cor:frame-dragging-tensor}, Lemma~\ref{lem:gps-redshift}): frame dragging is
a tensor-sector observable, and time dilation is the redshift law;
\item \textbf{Consequences} (Theorem~\ref{thm:gravity-consistent},
Corollary~\ref{cor:gravity-complete}): gravity is consistent with the canonical
architecture, complete within the hierarchy, with its single boundary disclosed.
\end{enumerate}

Gravity and spacetime are therefore emergent in The Actualization Theory: the gravitational
coupling is a read of the D96 spectral content, the tensor sector is the tensor face of
Difference against the reference $\eta$, the metric is the conformal construction from the
primitive pair, spacetime is emergent geometry with native dynamics, and the observables
--- from the Newton constant to the GPS offset --- are derived. The single Bekenstein
coefficient is a referenced boundary, not a derivation claim. The following chapter derives
the standard model from the same spectral foundation.

\begin{thebibliography}{99}
\bibitem{QG181} TQM-QG Phase 181, \emph{Newton Constant Origin}, Docs/Research.
\bibitem{QG182} TQM-QG Phase 182, \emph{G Bridge Origin}, Docs/Research.
\bibitem{QG183} TQM-QG Phase 183, \emph{Planck Scale Robustness}, Docs/Research.
\bibitem{QG184} TQM-QG Phase 184, \emph{Mass-Radius Origin}, Docs/Research.
\bibitem{QG186} TQM-QG Phase 186, \emph{Frame Dragging Origin}, Docs/Research.
\bibitem{QG187} TQM-QG Phase 187, \emph{GPS Correction Origin}, Docs/Research.
\bibitem{QG222} TQM-QG Phase 222, \emph{Native Metric Dynamics}, Docs/Research.
\bibitem{QG285} TQM-QG Phase 285, \emph{Psi Reinterpretation Audit}, Docs/Research.
\bibitem{QG286} TQM-QG Phase 286, \emph{Difference Duality Audit}, Docs/Research.
\bibitem{QG292} TQM-QG Phase 292, \emph{Foundation Stress Test}, Docs/Research.
\bibitem{MONO007} TQM-MONO007, \emph{Referee-Readiness Resolution}, Docs/Zenodo/Monograph.
\end{thebibliography}
 after the preamble
%         that defines the theorem environments (or include via the master file).
% ======================================================================

% ---- Theorem environments (define in the master preamble if not present) ----
% \newtheorem{definition}{Definition}[chapter]
% \newtheorem{lemma}[definition]{Lemma}
% \newtheorem{theorem}[definition]{Theorem}
% \newtheorem{corollary}[definition]{Corollary}
% \newtheorem{remark}[definition]{Remark}

\chapter{Gravity and Spacetime}
\label{chap:gravity}

\section{Introduction}
\label{sec:gravity-introduction}

The preceding chapters established the foundation of The Actualization Theory: the
primitives $\{\text{Difference},\eta\}$, the derived process of Actualization, the
inevitable D96 spectrum, the operator basis, and the emergent quantum structure. This
chapter derives \emph{gravity and spacetime} from the same foundation. We establish that
gravity is emergent --- a read of the counting measure and its tensor face --- and that
spacetime is emergent geometry \cite{QG181,QG184,QG222}.

The central results are: the gravitational coupling arises from the D96 spectral content
\cite{QG181,QG182,QG183}; the tensor sector is carried by the tensor face of Difference,
$\psi$, read against the tensor reference $\eta$ \cite{QG285,QG286}; the metric is
constructed from the counting measure and the reference \cite{QG292}; spacetime is emergent
geometry from the network and its dynamics \cite{QG222}; the black-hole relations and the
frame-dragging sector follow \cite{QG184,QG186}; and the GPS/gravitational-redshift
observables are the counting-measure redshift law \cite{QG187}.

Throughout, gravity is presented as \emph{emergent}, $\eta$ as the \emph{tensor
reference}, and $\psi$ as the \emph{tensor face of Difference}. The Bekenstein $1/4$
coefficient is \emph{excluded from the derivation claims} of this chapter and explicitly
referenced to the Boundary Layer, per the referee-safe architecture of MONO007
\cite{MONO007}.

We use only accepted canonical results and introduce no new physics and no speculation.

\section{Gravity from Difference}
\label{sec:gravity-from-difference}

\begin{definition}[Gravitational coupling]
\label{def:newton-constant}
The \emph{gravitational coupling} of the theory is the Newton constant $G$, derived from
the Planck scale
\begin{equation}
\label{eq:mpl}
M_{\mathrm{Pl}} \;=\; v\cdot(\Sigma m\cdot\#g\cdot\mathrm{occ}_2)^3,
\end{equation}
where $v$ is the weak scale, $\Sigma m = 95$ the total mode count, $\#g = 44$ the group
count, and $\mathrm{occ}_2$ the dense-band occupancy.
\end{definition}

\begin{theorem}[Gravity is derived from the D96 spectral content]
\label{thm:gravity-derived}
The gravitational coupling is derived from the D96 spectral content: the Planck scale is the
weak scale times the cube of the occupation-weighted spectral content, giving
$M_{\mathrm{Pl}} = 1.22335\times10^{19}\,\mathrm{GeV}$ and
$G = 6.6476\times10^{-11}\,\mathrm{m}^3\,\mathrm{kg}^{-1}\,\mathrm{s}^{-2}$.
\end{theorem}

\begin{proof}
The Newton-constant origin \cite{QG181} establishes the construction:
$M_{\mathrm{Pl}} = v\cdot(\Sigma m\cdot\#g\cdot\mathrm{occ}_2)^3 = 1.22335\times10^{19}$
GeV, matching the Planck scale within $0.2\%$, and $G = 1/M_{\mathrm{Pl}}^2$ within
$0.4\%$. The robustness audit \cite{QG183} confirms the physical exponent is cubic, and the
bridge audit \cite{QG182} shows the two G constructions are the same physical content.
Hence gravity is derived from the D96 spectral content --- the weak scale times the cube of
the occupation-weighted spectral structure.
\end{proof}

\begin{corollary}[Gravity is emergent]
\label{cor:gravity-emergent}
Gravity is emergent in the theory: the gravitational coupling is a derived function of the
D96 spectral content, not an independent input.
\end{corollary}

\begin{proof}
By Theorem~\ref{thm:gravity-derived}, $G$ is a function of the spectral constants
($\Sigma m$, $\#g$, $\mathrm{occ}_2$) and the weak scale. By Chapter~\ref{chap:d96}
(Theorem~\ref{thm:d96-derived}), the spectral constants are derived from the spectrum, which
is derived from the attractor. Hence gravity is twice-emergent: a read of the derived
spectral content.
\end{proof}

\section{The Tensor Sector}
\label{sec:tensor-sector}

\begin{theorem}[The tensor sector is the tensor face of Difference]
\label{thm:tensor-face}
The tensor sector of gravity is carried by the tensor face of Difference, $\psi$: the
Weyl content read against the tensor reference $\eta$.
\end{theorem}

\begin{proof}
The Difference duality \cite{QG286} establishes that $\rho$ and $\psi$ are the trace and
traceless faces of the one Difference (Chapter~\ref{chap:difference},
Theorem~\ref{thm:duality}). The psi-reinterpretation audit \cite{QG285} identifies $\psi$
as the Weyl content --- the difference from conformal flatness --- read against the
reference $\eta$ (Chapter~\ref{chap:eta}, Theorem~\ref{thm:tensor-read}). Hence the tensor
sector of gravity is carried by $\psi$, the tensor face of Difference, against the tensor
reference $\eta$.
\end{proof}

\begin{corollary}[The tensor sector is not a new primitive]
\label{cor:tensor-not-primitive}
The tensor sector introduces no new primitive: it is the tensor face of the existing
primitive Difference, read against the existing reference $\eta$.
\end{corollary}

\begin{proof}
By Theorem~\ref{thm:tensor-face}, the tensor content is $\psi$, a face of Difference, and
its reference is $\eta$. Both are members of the primitive pair established in
Chapter~\ref{chap:eta} (Theorem~\ref{thm:minimal-foundation}). Hence the tensor sector uses
no new primitive; it is the tensor face of the existing foundation.
\end{proof}

\begin{lemma}[The tensor observables require \texorpdfstring{$\eta$}{eta}]
\label{lem:tensor-requires-eta}
The tensor (gravitational) observables are defined against the tensor reference $\eta$:
the Weyl content, the metric completion, and the frame-dragging sector all presuppose the
reference.
\end{lemma}

\begin{proof}
The minimal-foundation result \cite{QG292} establishes that removing $\eta$ leaves the
scalar sector intact but destroys the tensor sub-sector. By Theorem~\ref{thm:tensor-face},
the tensor observables are reads of $\psi$ against $\eta$; hence they require $\eta$. This
is the canonical asymmetry: the scalar face needs no reference, the tensor face does.
\end{proof}

\section{The Metric Construction}
\label{sec:metric-construction}

\begin{definition}[The metric]
\label{def:metric}
The \emph{metric} of the theory is the conformal rescaling of the tensor reference:
\begin{equation}
\label{eq:metric}
g \;=\; \rho^{2/d}\,\eta,
\end{equation}
where $\rho$ is the scalar counting measure, $d$ the spatial dimension, and $\eta$ the
tensor reference.
\end{definition}

\begin{theorem}[The metric is constructed from the primitive pair]
\label{thm:metric-construction}
The metric is constructed from the primitive pair $\{\text{Difference},\eta\}$: the
conformal factor carries the scalar face of Difference ($\rho$), and the reference carries
the tensor structure ($\eta$).
\end{theorem}

\begin{proof}
The metric ansatz \eqref{eq:metric} is the canonical form established in
Chapter~\ref{chap:eta} (Theorem~\ref{thm:conformal}). The conformal factor $\rho^{2/d}$
carries the scalar counting content --- the scalar face of Difference --- and the reference
$\eta$ provides the tensor structure. Hence the metric is constructed from the primitive
pair, with no additional input. This construction is consistent with the minimal foundation
\cite{QG292}.
\end{proof}

\begin{corollary}[The metric is emergent]
\label{cor:metric-emergent}
The metric is an emergent construction: the conformal factor is derived from the counting
measure, and the reference is the primitive $\eta$; no independent metric postulate is
required.
\end{corollary}

\begin{proof}
By Theorem~\ref{thm:metric-construction}, the metric is built from $\rho$ and $\eta$. The
conformal factor is derived from the counting measure (the actualization share), and $\eta$
is the established primitive. Hence the metric is an emergent construction, not an imported
postulate.
\end{proof}

\section{Spacetime as Emergent Geometry}
\label{sec:emergent-geometry}

\begin{theorem}[Spacetime is emergent geometry]
\label{thm:spacetime-emergent}
Spacetime is emergent geometry in the theory: it is the geometry of the counting measure
and its dynamics, not a pre-existing arena.
\end{theorem}

\begin{proof}
The native-dynamics result \cite{QG222} establishes that the gravitational dynamics IS the
Q-event actualization flow: the branching process gives the counting measure $\rho_k =
\mu^k/S$ with count conservation by construction, and the branching continuity
$\rho_{k+1} = \mu\rho_k$ provides the native continuity/Noether statement. The metric is
constructed from $\rho$ and $\eta$ (Theorem~\ref{thm:metric-construction}). Hence spacetime
is the geometry of the counting measure --- emergent from the actualization flow, not a
pre-existing background.
\end{proof}

\begin{corollary}[No imported dynamics]
\label{cor:no-imported-dynamics}
The gravitational dynamics is native: it is the actualization flow, with no imported
Einstein or BDG dynamics.
\end{corollary}

\begin{proof}
By Theorem~\ref{thm:spacetime-emergent}, the dynamics is the Q-event actualization flow
\cite{QG222}. Hence no imported gravitational dynamics is required; the dynamics is derived
from the process itself.
\end{proof}

\section{Black Hole Relations}
\label{sec:black-hole}

\begin{theorem}[The mass-radius relation]
\label{thm:mass-radius}
The mass-radius relation $M \propto R$ follows from the counting structure: the per-octave
deficit profile gives $GM_{\mathrm{eff}} \propto R$.
\end{theorem}

\begin{proof}
The mass-radius origin \cite{QG184} establishes the relation: the deficit is per-octave/log
(the flat-rotation-curve profile), giving $a \propto -1/r$ and
$GM_{\mathrm{eff}} \propto R$. With the area-entropy relation $S \propto R^{d-1}$, the
Hawking temperature $T \propto 1/R$ is restored. Hence the black-hole relations follow from
the counting structure, with $M \propto R$ derived rather than assumed.
\end{proof}

\begin{theorem}[The Bekenstein coefficient is a boundary]
\label{thm:bekenstein-boundary}
The exact Bekenstein coefficient $S = A/4$ is \emph{not} a derivation claim of this
chapter: the structure $S \propto A$, $M \propto R$, $T \propto 1/R$ is derived, but the
exact $1/4$ coefficient requires an imported quantum factor and belongs to the Boundary
Layer.
\end{theorem}

\begin{proof}
The mass-radius and entropy relations are derived \cite{QG184}. The exact Bekenstein
$S = A/4$ coefficient, however, requires the imported $2\pi$ quantum factor
$T = \kappa/(2\pi)$, which the framework cannot produce internally; the coefficient is
therefore a documented boundary, not a derivation. Consistent with the referee-safe
architecture \cite{MONO007}, this chapter excludes the $1/4$ coefficient from its
derivation claims and references it to the Boundary Layer chapter of the monograph.
\end{proof}

\begin{corollary}[Boundary referenced, not claimed]
\label{cor:bekenstein-boundary-referenced}
The Bekenstein $1/4$ coefficient is referenced to the Boundary Layer and is not presented
as a derived result of the gravity chapter.
\end{corollary}

\begin{proof}
By Theorem~\ref{thm:bekenstein-boundary}, the exact coefficient is a documented boundary
requiring the imported quantum factor. Hence the gravity chapter presents only the derived
structure ($S \propto A$, $M \propto R$, $T \propto 1/R$) and references the coefficient to
the Boundary Layer, consistent with MONO007 \cite{MONO007}.
\end{proof}

\section{Frame Dragging}
\label{sec:frame-dragging}

\begin{definition}[Frame dragging]
\label{def:frame-dragging}
\emph{Frame dragging} is the gravitomagnetic effect by which a rotating source drags the
local inertial frames: the Lense--Thirring precession
\begin{equation}
\label{eq:frame-dragging}
\Omega_{\mathrm{LT}} \;=\; \frac{G\bigl(3(\mathbf{J}\cdot\hat{\mathbf{r}})\hat{\mathbf{r}}-\mathbf{J}\bigr)}{2c^2r^3}.
\end{equation}
\end{definition}

\begin{theorem}[Frame dragging is a tensor-sector observable]
\label{thm:frame-dragging}
Frame dragging is a $\psi$-sector observable: the gravitomagnetic $h_{0i}$ sector requires
the tensor face of Difference, and the Lense--Thirring precession follows.
\end{theorem}

\begin{proof}
The frame-dragging origin \cite{QG186} establishes the sector structure: the
conformally-flat $\rho$-only metric has $h_{0i} = 0$ (no frame dragging), while the tensor
face $\psi$ (spin-2) restores the full linearized Einstein sector including $h_{0i}$.
A rotating deficit field sources $\mathbf{J}$, giving the Lense--Thirring precession
\eqref{eq:frame-dragging}. The predicted precessions match observation: GP-B $41.1$ vs
$39.2$ mas/yr, LAGEOS $30.7$ vs $\sim31$ mas/yr, with the D96 $G$ shifting the results by
less than $1\%$. Hence frame dragging is a tensor-sector observable, carried by $\psi$
against $\eta$.
\end{proof}

\begin{corollary}[Frame dragging confirms the tensor face]
\label{cor:frame-dragging-tensor}
The frame-dragging sector confirms the role of $\psi$ as the tensor face of Difference:
without the tensor face there is no gravitomagnetic sector.
\end{corollary}

\begin{proof}
By Theorem~\ref{thm:frame-dragging}, the $h_{0i}$ sector vanishes in the $\rho$-only
metric and is restored by $\psi$. Hence frame dragging requires the tensor face of
Difference, confirming the duality structure \cite{QG286}.
\end{proof}

\begin{lemma}[Gravitational time dilation is the redshift law]
\label{lem:gps-redshift}
Gravitational time dilation is the counting-measure redshift law: the clock rate is
$d\tau/dt = \rho^{1/d} = \sqrt{-g_{00}}$, and the time dilation is the redshift
$\Delta\tau/\tau = (\rho_1/\rho_2)^{1/d} - 1$.
\end{lemma}

\begin{proof}
The GPS origin \cite{QG187} establishes that gravitational time dilation IS the redshift
law of the theory: the clock rate is $\rho^{1/d}$, the redshift is
$(\rho_1/\rho_2)^{1/d}-1$, and the weak-field limit gives the standard
$(GM/c^2)(1/r_1-1/r_2)$. The computed GPS offset $+38.5\ \mu$s/day matches the observed
$+38.6\ \mu$s/day within $0.2\%$. Hence time dilation is the counting-measure redshift law,
derived rather than imported.
\end{proof}

\section{Consequences}
\label{sec:gravity-consequences}

\begin{theorem}[Gravity is consistent with the canonical architecture]
\label{thm:gravity-consistent}
Gravity and spacetime are consistent with the canonical architecture: the gravitational
coupling is a read of the D96 spectral content, the tensor sector is the tensor face of
Difference against $\eta$, the metric is the conformal construction, spacetime is emergent
geometry, and the dynamics is the actualization flow.
\end{theorem}

\begin{proof}
The coupling: Theorem~\ref{thm:gravity-derived}; the tensor sector:
Theorem~\ref{thm:tensor-face}; the metric: Theorem~\ref{thm:metric-construction};
spacetime: Theorem~\ref{thm:spacetime-emergent}; the dynamics:
Theorem~\ref{thm:spacetime-emergent} (via \cite{QG222}). All are derived from the
foundation established in the preceding chapters, with no new primitives and no imported
dynamics.
\end{proof}

\begin{corollary}[The gravity pillar is complete within the hierarchy]
\label{cor:gravity-complete}
The gravity pillar of the theory is complete within the canonical hierarchy: every gravity
observable is a derived read of the D96 spectral content, the counting measure, and the
tensor face, with the Bekenstein $1/4$ coefficient the single referenced boundary.
\end{corollary}

\begin{proof}
By Theorem~\ref{thm:gravity-consistent}, all gravity content is derived; by
Theorem~\ref{thm:bekenstein-boundary}, the one exception (the exact $1/4$ coefficient) is
referenced to the Boundary Layer. Hence the gravity pillar is complete within the
hierarchy, with its single boundary explicitly disclosed.
\end{proof}

\begin{remark}[Referee-safe wording]
Consistent with MONO007 \cite{MONO007}, this chapter presents gravity as emergent, $\eta$ as
the tensor reference, and $\psi$ as the tensor face of Difference --- never as a new
primitive. The Bekenstein $1/4$ coefficient is excluded from the derivation claims and
referenced to the Boundary Layer. No claim here asserts more than the accepted results
establish.
\end{remark}

\section{Summary}
\label{sec:gravity-summary}

This chapter has derived gravity and spacetime. We have proved:

\begin{enumerate}
\item \textbf{Gravity from Difference} (Theorem~\ref{thm:gravity-derived},
Corollary~\ref{cor:gravity-emergent}): the gravitational coupling is derived from the D96
spectral content, and gravity is emergent;
\item \textbf{The tensor sector} (Theorem~\ref{thm:tensor-face},
Corollary~\ref{cor:tensor-not-primitive}, Lemma~\ref{lem:tensor-requires-eta}): the tensor
sector is the tensor face of Difference against $\eta$, not a new primitive;
\item \textbf{The metric construction} (Theorem~\ref{thm:metric-construction},
Corollary~\ref{cor:metric-emergent}): the metric is the conformal construction from the
primitive pair;
\item \textbf{Emergent geometry} (Theorem~\ref{thm:spacetime-emergent},
Corollary~\ref{cor:no-imported-dynamics}): spacetime is emergent geometry, with native
dynamics;
\item \textbf{Black hole relations} (Theorem~\ref{thm:mass-radius},
Theorem~\ref{thm:bekenstein-boundary}, Corollary~\ref{cor:bekenstein-boundary-referenced}):
$M \propto R$ is derived, and the Bekenstein $1/4$ coefficient is referenced to the
Boundary Layer;
\item \textbf{Frame dragging} (Theorem~\ref{thm:frame-dragging},
Corollary~\ref{cor:frame-dragging-tensor}, Lemma~\ref{lem:gps-redshift}): frame dragging is
a tensor-sector observable, and time dilation is the redshift law;
\item \textbf{Consequences} (Theorem~\ref{thm:gravity-consistent},
Corollary~\ref{cor:gravity-complete}): gravity is consistent with the canonical
architecture, complete within the hierarchy, with its single boundary disclosed.
\end{enumerate}

Gravity and spacetime are therefore emergent in The Actualization Theory: the gravitational
coupling is a read of the D96 spectral content, the tensor sector is the tensor face of
Difference against the reference $\eta$, the metric is the conformal construction from the
primitive pair, spacetime is emergent geometry with native dynamics, and the observables
--- from the Newton constant to the GPS offset --- are derived. The single Bekenstein
coefficient is a referenced boundary, not a derivation claim. The following chapter derives
the standard model from the same spectral foundation.

\begin{thebibliography}{99}
\bibitem{QG181} TQM-QG Phase 181, \emph{Newton Constant Origin}, Docs/Research.
\bibitem{QG182} TQM-QG Phase 182, \emph{G Bridge Origin}, Docs/Research.
\bibitem{QG183} TQM-QG Phase 183, \emph{Planck Scale Robustness}, Docs/Research.
\bibitem{QG184} TQM-QG Phase 184, \emph{Mass-Radius Origin}, Docs/Research.
\bibitem{QG186} TQM-QG Phase 186, \emph{Frame Dragging Origin}, Docs/Research.
\bibitem{QG187} TQM-QG Phase 187, \emph{GPS Correction Origin}, Docs/Research.
\bibitem{QG222} TQM-QG Phase 222, \emph{Native Metric Dynamics}, Docs/Research.
\bibitem{QG285} TQM-QG Phase 285, \emph{Psi Reinterpretation Audit}, Docs/Research.
\bibitem{QG286} TQM-QG Phase 286, \emph{Difference Duality Audit}, Docs/Research.
\bibitem{QG292} TQM-QG Phase 292, \emph{Foundation Stress Test}, Docs/Research.
\bibitem{MONO007} TQM-MONO007, \emph{Referee-Readiness Resolution}, Docs/Zenodo/Monograph.
\end{thebibliography}
 after the preamble
%         that defines the theorem environments (or include via the master file).
% ======================================================================

% ---- Theorem environments (define in the master preamble if not present) ----
% \newtheorem{definition}{Definition}[chapter]
% \newtheorem{lemma}[definition]{Lemma}
% \newtheorem{theorem}[definition]{Theorem}
% \newtheorem{corollary}[definition]{Corollary}
% \newtheorem{remark}[definition]{Remark}

\chapter{Gravity and Spacetime}
\label{chap:gravity}

\section{Introduction}
\label{sec:gravity-introduction}

The preceding chapters established the foundation of The Actualization Theory: the
primitives $\{\text{Difference},\eta\}$, the derived process of Actualization, the
inevitable D96 spectrum, the operator basis, and the emergent quantum structure. This
chapter derives \emph{gravity and spacetime} from the same foundation. We establish that
gravity is emergent --- a read of the counting measure and its tensor face --- and that
spacetime is emergent geometry \cite{QG181,QG184,QG222}.

The central results are: the gravitational coupling arises from the D96 spectral content
\cite{QG181,QG182,QG183}; the tensor sector is carried by the tensor face of Difference,
$\psi$, read against the tensor reference $\eta$ \cite{QG285,QG286}; the metric is
constructed from the counting measure and the reference \cite{QG292}; spacetime is emergent
geometry from the network and its dynamics \cite{QG222}; the black-hole relations and the
frame-dragging sector follow \cite{QG184,QG186}; and the GPS/gravitational-redshift
observables are the counting-measure redshift law \cite{QG187}.

Throughout, gravity is presented as \emph{emergent}, $\eta$ as the \emph{tensor
reference}, and $\psi$ as the \emph{tensor face of Difference}. The Bekenstein $1/4$
coefficient is \emph{excluded from the derivation claims} of this chapter and explicitly
referenced to the Boundary Layer, per the referee-safe architecture of MONO007
\cite{MONO007}.

We use only accepted canonical results and introduce no new physics and no speculation.

\section{Gravity from Difference}
\label{sec:gravity-from-difference}

\begin{definition}[Gravitational coupling]
\label{def:newton-constant}
The \emph{gravitational coupling} of the theory is the Newton constant $G$, derived from
the Planck scale
\begin{equation}
\label{eq:mpl}
M_{\mathrm{Pl}} \;=\; v\cdot(\Sigma m\cdot\#g\cdot\mathrm{occ}_2)^3,
\end{equation}
where $v$ is the weak scale, $\Sigma m = 95$ the total mode count, $\#g = 44$ the group
count, and $\mathrm{occ}_2$ the dense-band occupancy.
\end{definition}

\begin{theorem}[Gravity is derived from the D96 spectral content]
\label{thm:gravity-derived}
The gravitational coupling is derived from the D96 spectral content: the Planck scale is the
weak scale times the cube of the occupation-weighted spectral content, giving
$M_{\mathrm{Pl}} = 1.22335\times10^{19}\,\mathrm{GeV}$ and
$G = 6.6476\times10^{-11}\,\mathrm{m}^3\,\mathrm{kg}^{-1}\,\mathrm{s}^{-2}$.
\end{theorem}

\begin{proof}
The Newton-constant origin \cite{QG181} establishes the construction:
$M_{\mathrm{Pl}} = v\cdot(\Sigma m\cdot\#g\cdot\mathrm{occ}_2)^3 = 1.22335\times10^{19}$
GeV, matching the Planck scale within $0.2\%$, and $G = 1/M_{\mathrm{Pl}}^2$ within
$0.4\%$. The robustness audit \cite{QG183} confirms the physical exponent is cubic, and the
bridge audit \cite{QG182} shows the two G constructions are the same physical content.
Hence gravity is derived from the D96 spectral content --- the weak scale times the cube of
the occupation-weighted spectral structure.
\end{proof}

\begin{corollary}[Gravity is emergent]
\label{cor:gravity-emergent}
Gravity is emergent in the theory: the gravitational coupling is a derived function of the
D96 spectral content, not an independent input.
\end{corollary}

\begin{proof}
By Theorem~\ref{thm:gravity-derived}, $G$ is a function of the spectral constants
($\Sigma m$, $\#g$, $\mathrm{occ}_2$) and the weak scale. By Chapter~\ref{chap:d96}
(Theorem~\ref{thm:d96-derived}), the spectral constants are derived from the spectrum, which
is derived from the attractor. Hence gravity is twice-emergent: a read of the derived
spectral content.
\end{proof}

\section{The Tensor Sector}
\label{sec:tensor-sector}

\begin{theorem}[The tensor sector is the tensor face of Difference]
\label{thm:tensor-face}
The tensor sector of gravity is carried by the tensor face of Difference, $\psi$: the
Weyl content read against the tensor reference $\eta$.
\end{theorem}

\begin{proof}
The Difference duality \cite{QG286} establishes that $\rho$ and $\psi$ are the trace and
traceless faces of the one Difference (Chapter~\ref{chap:difference},
Theorem~\ref{thm:duality}). The psi-reinterpretation audit \cite{QG285} identifies $\psi$
as the Weyl content --- the difference from conformal flatness --- read against the
reference $\eta$ (Chapter~\ref{chap:eta}, Theorem~\ref{thm:tensor-read}). Hence the tensor
sector of gravity is carried by $\psi$, the tensor face of Difference, against the tensor
reference $\eta$.
\end{proof}

\begin{corollary}[The tensor sector is not a new primitive]
\label{cor:tensor-not-primitive}
The tensor sector introduces no new primitive: it is the tensor face of the existing
primitive Difference, read against the existing reference $\eta$.
\end{corollary}

\begin{proof}
By Theorem~\ref{thm:tensor-face}, the tensor content is $\psi$, a face of Difference, and
its reference is $\eta$. Both are members of the primitive pair established in
Chapter~\ref{chap:eta} (Theorem~\ref{thm:minimal-foundation}). Hence the tensor sector uses
no new primitive; it is the tensor face of the existing foundation.
\end{proof}

\begin{lemma}[The tensor observables require \texorpdfstring{$\eta$}{eta}]
\label{lem:tensor-requires-eta}
The tensor (gravitational) observables are defined against the tensor reference $\eta$:
the Weyl content, the metric completion, and the frame-dragging sector all presuppose the
reference.
\end{lemma}

\begin{proof}
The minimal-foundation result \cite{QG292} establishes that removing $\eta$ leaves the
scalar sector intact but destroys the tensor sub-sector. By Theorem~\ref{thm:tensor-face},
the tensor observables are reads of $\psi$ against $\eta$; hence they require $\eta$. This
is the canonical asymmetry: the scalar face needs no reference, the tensor face does.
\end{proof}

\section{The Metric Construction}
\label{sec:metric-construction}

\begin{definition}[The metric]
\label{def:metric}
The \emph{metric} of the theory is the conformal rescaling of the tensor reference:
\begin{equation}
\label{eq:metric}
g \;=\; \rho^{2/d}\,\eta,
\end{equation}
where $\rho$ is the scalar counting measure, $d$ the spatial dimension, and $\eta$ the
tensor reference.
\end{definition}

\begin{theorem}[The metric is constructed from the primitive pair]
\label{thm:metric-construction}
The metric is constructed from the primitive pair $\{\text{Difference},\eta\}$: the
conformal factor carries the scalar face of Difference ($\rho$), and the reference carries
the tensor structure ($\eta$).
\end{theorem}

\begin{proof}
The metric ansatz \eqref{eq:metric} is the canonical form established in
Chapter~\ref{chap:eta} (Theorem~\ref{thm:conformal}). The conformal factor $\rho^{2/d}$
carries the scalar counting content --- the scalar face of Difference --- and the reference
$\eta$ provides the tensor structure. Hence the metric is constructed from the primitive
pair, with no additional input. This construction is consistent with the minimal foundation
\cite{QG292}.
\end{proof}

\begin{corollary}[The metric is emergent]
\label{cor:metric-emergent}
The metric is an emergent construction: the conformal factor is derived from the counting
measure, and the reference is the primitive $\eta$; no independent metric postulate is
required.
\end{corollary}

\begin{proof}
By Theorem~\ref{thm:metric-construction}, the metric is built from $\rho$ and $\eta$. The
conformal factor is derived from the counting measure (the actualization share), and $\eta$
is the established primitive. Hence the metric is an emergent construction, not an imported
postulate.
\end{proof}

\section{Spacetime as Emergent Geometry}
\label{sec:emergent-geometry}

\begin{theorem}[Spacetime is emergent geometry]
\label{thm:spacetime-emergent}
Spacetime is emergent geometry in the theory: it is the geometry of the counting measure
and its dynamics, not a pre-existing arena.
\end{theorem}

\begin{proof}
The native-dynamics result \cite{QG222} establishes that the gravitational dynamics IS the
Q-event actualization flow: the branching process gives the counting measure $\rho_k =
\mu^k/S$ with count conservation by construction, and the branching continuity
$\rho_{k+1} = \mu\rho_k$ provides the native continuity/Noether statement. The metric is
constructed from $\rho$ and $\eta$ (Theorem~\ref{thm:metric-construction}). Hence spacetime
is the geometry of the counting measure --- emergent from the actualization flow, not a
pre-existing background.
\end{proof}

\begin{corollary}[No imported dynamics]
\label{cor:no-imported-dynamics}
The gravitational dynamics is native: it is the actualization flow, with no imported
Einstein or BDG dynamics.
\end{corollary}

\begin{proof}
By Theorem~\ref{thm:spacetime-emergent}, the dynamics is the Q-event actualization flow
\cite{QG222}. Hence no imported gravitational dynamics is required; the dynamics is derived
from the process itself.
\end{proof}

\section{Black Hole Relations}
\label{sec:black-hole}

\begin{theorem}[The mass-radius relation]
\label{thm:mass-radius}
The mass-radius relation $M \propto R$ follows from the counting structure: the per-octave
deficit profile gives $GM_{\mathrm{eff}} \propto R$.
\end{theorem}

\begin{proof}
The mass-radius origin \cite{QG184} establishes the relation: the deficit is per-octave/log
(the flat-rotation-curve profile), giving $a \propto -1/r$ and
$GM_{\mathrm{eff}} \propto R$. With the area-entropy relation $S \propto R^{d-1}$, the
Hawking temperature $T \propto 1/R$ is restored. Hence the black-hole relations follow from
the counting structure, with $M \propto R$ derived rather than assumed.
\end{proof}

\begin{theorem}[The Bekenstein coefficient is a boundary]
\label{thm:bekenstein-boundary}
The exact Bekenstein coefficient $S = A/4$ is \emph{not} a derivation claim of this
chapter: the structure $S \propto A$, $M \propto R$, $T \propto 1/R$ is derived, but the
exact $1/4$ coefficient requires an imported quantum factor and belongs to the Boundary
Layer.
\end{theorem}

\begin{proof}
The mass-radius and entropy relations are derived \cite{QG184}. The exact Bekenstein
$S = A/4$ coefficient, however, requires the imported $2\pi$ quantum factor
$T = \kappa/(2\pi)$, which the framework cannot produce internally; the coefficient is
therefore a documented boundary, not a derivation. Consistent with the referee-safe
architecture \cite{MONO007}, this chapter excludes the $1/4$ coefficient from its
derivation claims and references it to the Boundary Layer chapter of the monograph.
\end{proof}

\begin{corollary}[Boundary referenced, not claimed]
\label{cor:bekenstein-boundary-referenced}
The Bekenstein $1/4$ coefficient is referenced to the Boundary Layer and is not presented
as a derived result of the gravity chapter.
\end{corollary}

\begin{proof}
By Theorem~\ref{thm:bekenstein-boundary}, the exact coefficient is a documented boundary
requiring the imported quantum factor. Hence the gravity chapter presents only the derived
structure ($S \propto A$, $M \propto R$, $T \propto 1/R$) and references the coefficient to
the Boundary Layer, consistent with MONO007 \cite{MONO007}.
\end{proof}

\section{Frame Dragging}
\label{sec:frame-dragging}

\begin{definition}[Frame dragging]
\label{def:frame-dragging}
\emph{Frame dragging} is the gravitomagnetic effect by which a rotating source drags the
local inertial frames: the Lense--Thirring precession
\begin{equation}
\label{eq:frame-dragging}
\Omega_{\mathrm{LT}} \;=\; \frac{G\bigl(3(\mathbf{J}\cdot\hat{\mathbf{r}})\hat{\mathbf{r}}-\mathbf{J}\bigr)}{2c^2r^3}.
\end{equation}
\end{definition}

\begin{theorem}[Frame dragging is a tensor-sector observable]
\label{thm:frame-dragging}
Frame dragging is a $\psi$-sector observable: the gravitomagnetic $h_{0i}$ sector requires
the tensor face of Difference, and the Lense--Thirring precession follows.
\end{theorem}

\begin{proof}
The frame-dragging origin \cite{QG186} establishes the sector structure: the
conformally-flat $\rho$-only metric has $h_{0i} = 0$ (no frame dragging), while the tensor
face $\psi$ (spin-2) restores the full linearized Einstein sector including $h_{0i}$.
A rotating deficit field sources $\mathbf{J}$, giving the Lense--Thirring precession
\eqref{eq:frame-dragging}. The predicted precessions match observation: GP-B $41.1$ vs
$39.2$ mas/yr, LAGEOS $30.7$ vs $\sim31$ mas/yr, with the D96 $G$ shifting the results by
less than $1\%$. Hence frame dragging is a tensor-sector observable, carried by $\psi$
against $\eta$.
\end{proof}

\begin{corollary}[Frame dragging confirms the tensor face]
\label{cor:frame-dragging-tensor}
The frame-dragging sector confirms the role of $\psi$ as the tensor face of Difference:
without the tensor face there is no gravitomagnetic sector.
\end{corollary}

\begin{proof}
By Theorem~\ref{thm:frame-dragging}, the $h_{0i}$ sector vanishes in the $\rho$-only
metric and is restored by $\psi$. Hence frame dragging requires the tensor face of
Difference, confirming the duality structure \cite{QG286}.
\end{proof}

\begin{lemma}[Gravitational time dilation is the redshift law]
\label{lem:gps-redshift}
Gravitational time dilation is the counting-measure redshift law: the clock rate is
$d\tau/dt = \rho^{1/d} = \sqrt{-g_{00}}$, and the time dilation is the redshift
$\Delta\tau/\tau = (\rho_1/\rho_2)^{1/d} - 1$.
\end{lemma}

\begin{proof}
The GPS origin \cite{QG187} establishes that gravitational time dilation IS the redshift
law of the theory: the clock rate is $\rho^{1/d}$, the redshift is
$(\rho_1/\rho_2)^{1/d}-1$, and the weak-field limit gives the standard
$(GM/c^2)(1/r_1-1/r_2)$. The computed GPS offset $+38.5\ \mu$s/day matches the observed
$+38.6\ \mu$s/day within $0.2\%$. Hence time dilation is the counting-measure redshift law,
derived rather than imported.
\end{proof}

\section{Consequences}
\label{sec:gravity-consequences}

\begin{theorem}[Gravity is consistent with the canonical architecture]
\label{thm:gravity-consistent}
Gravity and spacetime are consistent with the canonical architecture: the gravitational
coupling is a read of the D96 spectral content, the tensor sector is the tensor face of
Difference against $\eta$, the metric is the conformal construction, spacetime is emergent
geometry, and the dynamics is the actualization flow.
\end{theorem}

\begin{proof}
The coupling: Theorem~\ref{thm:gravity-derived}; the tensor sector:
Theorem~\ref{thm:tensor-face}; the metric: Theorem~\ref{thm:metric-construction};
spacetime: Theorem~\ref{thm:spacetime-emergent}; the dynamics:
Theorem~\ref{thm:spacetime-emergent} (via \cite{QG222}). All are derived from the
foundation established in the preceding chapters, with no new primitives and no imported
dynamics.
\end{proof}

\begin{corollary}[The gravity pillar is complete within the hierarchy]
\label{cor:gravity-complete}
The gravity pillar of the theory is complete within the canonical hierarchy: every gravity
observable is a derived read of the D96 spectral content, the counting measure, and the
tensor face, with the Bekenstein $1/4$ coefficient the single referenced boundary.
\end{corollary}

\begin{proof}
By Theorem~\ref{thm:gravity-consistent}, all gravity content is derived; by
Theorem~\ref{thm:bekenstein-boundary}, the one exception (the exact $1/4$ coefficient) is
referenced to the Boundary Layer. Hence the gravity pillar is complete within the
hierarchy, with its single boundary explicitly disclosed.
\end{proof}

\begin{remark}[Referee-safe wording]
Consistent with MONO007 \cite{MONO007}, this chapter presents gravity as emergent, $\eta$ as
the tensor reference, and $\psi$ as the tensor face of Difference --- never as a new
primitive. The Bekenstein $1/4$ coefficient is excluded from the derivation claims and
referenced to the Boundary Layer. No claim here asserts more than the accepted results
establish.
\end{remark}

\section{Summary}
\label{sec:gravity-summary}

This chapter has derived gravity and spacetime. We have proved:

\begin{enumerate}
\item \textbf{Gravity from Difference} (Theorem~\ref{thm:gravity-derived},
Corollary~\ref{cor:gravity-emergent}): the gravitational coupling is derived from the D96
spectral content, and gravity is emergent;
\item \textbf{The tensor sector} (Theorem~\ref{thm:tensor-face},
Corollary~\ref{cor:tensor-not-primitive}, Lemma~\ref{lem:tensor-requires-eta}): the tensor
sector is the tensor face of Difference against $\eta$, not a new primitive;
\item \textbf{The metric construction} (Theorem~\ref{thm:metric-construction},
Corollary~\ref{cor:metric-emergent}): the metric is the conformal construction from the
primitive pair;
\item \textbf{Emergent geometry} (Theorem~\ref{thm:spacetime-emergent},
Corollary~\ref{cor:no-imported-dynamics}): spacetime is emergent geometry, with native
dynamics;
\item \textbf{Black hole relations} (Theorem~\ref{thm:mass-radius},
Theorem~\ref{thm:bekenstein-boundary}, Corollary~\ref{cor:bekenstein-boundary-referenced}):
$M \propto R$ is derived, and the Bekenstein $1/4$ coefficient is referenced to the
Boundary Layer;
\item \textbf{Frame dragging} (Theorem~\ref{thm:frame-dragging},
Corollary~\ref{cor:frame-dragging-tensor}, Lemma~\ref{lem:gps-redshift}): frame dragging is
a tensor-sector observable, and time dilation is the redshift law;
\item \textbf{Consequences} (Theorem~\ref{thm:gravity-consistent},
Corollary~\ref{cor:gravity-complete}): gravity is consistent with the canonical
architecture, complete within the hierarchy, with its single boundary disclosed.
\end{enumerate}

Gravity and spacetime are therefore emergent in The Actualization Theory: the gravitational
coupling is a read of the D96 spectral content, the tensor sector is the tensor face of
Difference against the reference $\eta$, the metric is the conformal construction from the
primitive pair, spacetime is emergent geometry with native dynamics, and the observables
--- from the Newton constant to the GPS offset --- are derived. The single Bekenstein
coefficient is a referenced boundary, not a derivation claim. The following chapter derives
the standard model from the same spectral foundation.

\begin{thebibliography}{99}
\bibitem{QG181} TQM-QG Phase 181, \emph{Newton Constant Origin}, Docs/Research.
\bibitem{QG182} TQM-QG Phase 182, \emph{G Bridge Origin}, Docs/Research.
\bibitem{QG183} TQM-QG Phase 183, \emph{Planck Scale Robustness}, Docs/Research.
\bibitem{QG184} TQM-QG Phase 184, \emph{Mass-Radius Origin}, Docs/Research.
\bibitem{QG186} TQM-QG Phase 186, \emph{Frame Dragging Origin}, Docs/Research.
\bibitem{QG187} TQM-QG Phase 187, \emph{GPS Correction Origin}, Docs/Research.
\bibitem{QG222} TQM-QG Phase 222, \emph{Native Metric Dynamics}, Docs/Research.
\bibitem{QG285} TQM-QG Phase 285, \emph{Psi Reinterpretation Audit}, Docs/Research.
\bibitem{QG286} TQM-QG Phase 286, \emph{Difference Duality Audit}, Docs/Research.
\bibitem{QG292} TQM-QG Phase 292, \emph{Foundation Stress Test}, Docs/Research.
\bibitem{MONO007} TQM-MONO007, \emph{Referee-Readiness Resolution}, Docs/Zenodo/Monograph.
\end{thebibliography}
